\chapter{System Design}

The system design chapter translates the requirements identified in the previous chapter into a comprehensive technical blueprint for the AI-Integrated Washing Machine. This chapter covers the system architecture, data flow diagrams, hardware design, software architecture, database design, module specifications, and integration details.

\section{System Architecture}

The AI-Integrated Washing Machine follows a four-layer architecture that separates concerns between hardware interaction, processing logic, intelligent decision-making, and user communication. This layered approach ensures modularity, maintainability, and scalability.

\begin{figure}[H]
\centering
\fbox{\begin{minipage}{0.9\textwidth}
\centering
\textbf{Four-Layer Architecture}
\begin{itemize}
    \item \textbf{User Interaction Layer:} Touchscreen display, optional mobile app, physical buttons (manual override)
    \item \textbf{Software Layer:} Fabric classification models (CNN), cycle generation algorithms, control logic \& state machine, UI framework
    \item \textbf{Processing Layer:} Raspberry Pi (image processing, ML inference, main app), Arduino/microcontroller (sensor data acquisition, actuator control)
    \item \textbf{Physical Layer:} Camera module, load cells, turbidity sensor, temperature sensor, water valves, detergent dispenser, drum motor, drain pump
\end{itemize}
\end{minipage}}
\caption{AI-Integrated Washing Machine Four-Layer Architecture}
\label{fig:fourlayer}
\end{figure}

\subsection{Layer Descriptions}

\begin{itemize}
    \item \textbf{Physical Layer:} Contains all hardware components including sensors (camera, load cells, turbidity sensor, temperature sensor) and actuators (water valves, detergent dispensers, drum motor, drain pump). This layer interfaces directly with the physical laundry environment.
    \item \textbf{Processing Layer:} Comprises the embedded computing hardware. A Raspberry Pi handles high-level processing including image capture and machine learning inference. An Arduino or similar microcontroller manages real-time sensor data acquisition and precise actuator control.
    \item \textbf{Software Layer:} Contains all algorithms and logic including trained CNN models for fabric classification, cycle generation algorithms, control state machine, and user interface software.
    \item \textbf{User Interaction Layer:} Provides interfaces for user communication including the touchscreen display, optional mobile application connectivity, and physical buttons for manual override.
\end{itemize}

\section{Data Flow Diagrams}

\subsection{Level 0 DFD (Context Diagram)}

\begin{figure}[H]
\centering
\fbox{\begin{minipage}{0.8\textwidth}
\centering
\textbf{Context Diagram (Level 0 DFD)}
\begin{itemize}
    \item External entity: User
    \item Inputs: Load clothes, manual input
    \item System: AI-Integrated Washing Machine (contains Fabric Data Storage, Wash Cycle Generation, Hardware Components, Control System)
    \item Outputs: Display feedback
\end{itemize}
\end{minipage}}
\caption{Context Diagram (Level 0 DFD)}
\label{fig:dfd0}
\end{figure}

\subsection{Level 1 DFD}

\begin{figure}[H]
\centering
\fbox{\begin{minipage}{0.9\textwidth}
\centering
\textbf{Level 1 Data Flow Diagram}
\begin{itemize}
    \item User $\rightarrow$ Camera Module: Load clothes
    \item Camera Module $\rightarrow$ Vision Processing: Images
    \item Vision Processing $\rightarrow$ Cycle Generation: Fabric type
    \item Sensors (Load, Turbidity, Temp) $\rightarrow$ Data Logger $\rightarrow$ Cycle Generation: Sensor data
    \item Cycle Generation $\rightarrow$ Control System: Cycle parameters
    \item Control System $\rightarrow$ Actuators (valves, dispenser, motor): Control signals
    \item Control System $\rightarrow$ Touchscreen Display / Mobile App: Status updates
    \item User $\rightarrow$ Touchscreen / Mobile App: Manual override, feedback
\end{itemize}
\end{minipage}}
\caption{Level 1 Data Flow Diagram}
\label{fig:dfd1}
\end{figure}

\section{Hardware Design}

\subsection{Component Selection}

\begin{table}[H]
\centering
\caption{Hardware Component Specifications}
\label{tab:hardware}
\begin{tabular}{|l|p{3.5cm}|p{6cm}|}
\hline
\textbf{Component} & \textbf{Specification} & \textbf{Purpose} \\
\hline
Camera Module & Raspberry Pi Camera Module 3 (12MP) or USB camera with autofocus & Capture high-resolution images of clothing for fabric analysis \\
\hline
Main Processor & Raspberry Pi 4 (4GB RAM) & Run ML models, image processing, main application logic \\
\hline
Microcontroller & Arduino Uno / ESP32 & Real-time sensor reading and actuator control \\
\hline
Load Cells & 4x 50kg load cells with HX711 amplifier & Measure total laundry weight \\
\hline
Turbidity Sensor & Analog turbidity sensor & Detect water clarity for soil level assessment \\
\hline
Temperature Sensor & DS18B20 waterproof sensor & Monitor water temperature \\
\hline
Water Valves & 12V solenoid valves (hot/cold) & Control water inflow \\
\hline
Detergent Dispenser & Peristaltic pump with stepper motor & Precise liquid detergent dispensing \\
\hline
Drum Motor & BLDC motor with controller & Variable speed drum control \\
\hline
Display & 7-inch touchscreen display & User interface \\
\hline
LED Lighting & High-CRI white LEDs & Consistent illumination for camera \\
\hline
\end{tabular}
\end{table}

\subsection{Hardware Integration Diagram}

\begin{figure}[H]
\centering
\fbox{\begin{minipage}{0.9\textwidth}
\centering
\textbf{Hardware Integration}
\begin{itemize}
    \item Power Supply (230V AC $\rightarrow$ 12V/5V) supplies Raspberry Pi and other components.
    \item Raspberry Pi connects to Camera Module, Touchscreen Display, and Arduino (via USB).
    \item Arduino interfaces with: Load Cells (x4 + HX711), Turbidity Sensor, Temperature Sensor, Motor Driver.
    \item Motor Driver controls Drum Motor (BLDC).
    \item Arduino also controls Water Valves (hot/cold), Detergent Pump (peristaltic), and Drain Pump.
\end{itemize}
\end{minipage}}
\caption{Hardware Integration Diagram}
\label{fig:hw}
\end{figure}

\section{Software Design}

\subsection{Software Architecture}

The software architecture follows a modular design with clear separation between components:

\begin{figure}[H]
\centering
\fbox{\begin{minipage}{0.9\textwidth}
\centering
\textbf{Software Architecture}
\begin{itemize}
    \item \textbf{Main Application} (Python/Node.js/C++) orchestrates all modules.
    \item Modules:
        \begin{itemize}
            \item \textbf{Vision Module}: Image capture, preprocessing, inference, fabric detection.
            \item \textbf{Sensor Module}: Read load cells, turbidity, temperature.
            \item \textbf{Control Module}: State machine, cycle execution, safety monitoring.
            \item \textbf{Cycle Generator}: Parameter calculation, resource optimization, cycle selection.
            \item \textbf{User Interface}: Display updates, touch handling, mobile API.
            \item \textbf{Actuator Control}: Valve, dispenser, motor control.
        \end{itemize}
    \item \textbf{Database Manager}: SQLite / local storage for fabric knowledge, wash history, user preferences.
\end{itemize}
\end{minipage}}
\caption{Software Architecture Diagram}
\label{fig:swarch}
\end{figure}

\subsection{Fabric Classification Model}

The fabric detection system uses a Convolutional Neural Network (CNN) trained on a dataset of fabric images. The model architecture is based on MobileNetV2 for efficient embedded deployment.

\begin{verbatim}
--------------------------------------------------
Input Layer: 224x224x3 RGB Image
|
v
MobileNetV2 Base (Pretrained on ImageNet)
|
v
Global Average Pooling
|
v
Dense Layer (128 units, ReLU activation)
|
v
Dropout (0.3)
|
v
Output Layer (Softmax, 8 fabric classes)
   - Cotton         - Polyester
   - Wool           - Silk
   - Linen          - Nylon
   - Rayon          - Blends
--------------------------------------------------
Training Configuration:
- Dataset: 10,000+ labeled fabric images
- Split: 70% training, 15% validation, 15% testing
- Optimizer: Adam (learning rate = 0.001)
- Loss Function: Categorical Crossentropy
- Batch Size: 32
- Epochs: 50 (with early stopping)
- Data Augmentation: Rotation, zoom, brightness, contrast
\end{verbatim}

\section{Database Design}

\subsection{Database Schema}

The system uses a lightweight SQLite database for local storage of fabric knowledge, wash history, and user preferences.

\begin{figure}[H]
\centering
\fbox{\begin{minipage}{0.9\textwidth}
\centering
\textbf{Database Schema}
\begin{itemize}
    \item \texttt{fabric\_types} (id, name, category, delicate\_flag, default\_temp, max\_temp, agitation\_speed, spin\_speed, notes)
    \item \texttt{wash\_cycles} (id, fabric\_type\_id, water\_temp, water\_level, detergent\_amount, duration, spin\_speed, rinse\_count)
    \item \texttt{fabric\_knowledge} (id, fabric\_type\_id, weight\_factor, soil\_factor, detergent\_factor, special\_handling, last\_updated)
    \item \texttt{wash\_history} (id, fabric\_type\_id, timestamp, load\_weight, soil\_level, cycle\_id, water\_used, detergent\_used, cycle\_complete, user\_rating)
    \item \texttt{cycle\_log} (id, wash\_id, phase, start\_time, end\_time, sensor\_readings)
\end{itemize}
Relationships: foreign keys link tables appropriately.
\end{minipage}}
\caption{Database Schema Diagram}
\label{fig:db}
\end{figure}

\subsection{Table Definitions}

\begin{lstlisting}[language=SQL, caption={fabric\_types table}, label={lst:fabric}]
CREATE TABLE fabric_types (
    id INTEGER PRIMARY KEY AUTOINCREMENT,
    name VARCHAR(50) NOT NULL,
    category VARCHAR(30),
    delicate_flag BOOLEAN DEFAULT 0,
    default_temp INTEGER,
    max_temp INTEGER,
    agitation_speed VARCHAR(20),
    spin_speed INTEGER,
    notes TEXT,
    created_at TIMESTAMP DEFAULT CURRENT_TIMESTAMP
);
CREATE INDEX idx_fabric_types_name ON fabric_types(name);
\end{lstlisting}

\begin{lstlisting}[language=SQL, caption={wash\_cycles table}, label={lst:cycles}]
CREATE TABLE wash_cycles (
    id INTEGER PRIMARY KEY AUTOINCREMENT,
    fabric_type_id INTEGER,
    water_temp INTEGER,
    water_level VARCHAR(20),
    detergent_amount DECIMAL(5,2),
    duration INTEGER,
    spin_speed INTEGER,
    rinse_count INTEGER DEFAULT 2,
    FOREIGN KEY (fabric_type_id) REFERENCES fabric_types(id)
);
CREATE INDEX idx_wash_cycles_fabric ON wash_cycles(fabric_type_id);
\end{lstlisting}

\begin{lstlisting}[language=SQL, caption={wash\_history table}, label={lst:history}]
CREATE TABLE wash_history (
    id INTEGER PRIMARY KEY AUTOINCREMENT,
    fabric_type_id INTEGER,
    timestamp DATETIME DEFAULT CURRENT_TIMESTAMP,
    load_weight DECIMAL(5,2),
    soil_level INTEGER,
    cycle_id INTEGER,
    water_used DECIMAL(6,2),
    detergent_used DECIMAL(5,2),
    cycle_complete BOOLEAN DEFAULT 0,
    user_rating INTEGER,
    FOREIGN KEY (fabric_type_id) REFERENCES fabric_types(id),
    FOREIGN KEY (cycle_id) REFERENCES wash_cycles(id)
);
CREATE INDEX idx_wash_history_timestamp ON wash_history(timestamp);
CREATE INDEX idx_wash_history_fabric ON wash_history(fabric_type_id);
\end{lstlisting}

\begin{lstlisting}[language=SQL, caption={cycle\_log table}, label={lst:log}]
CREATE TABLE cycle_log (
    id INTEGER PRIMARY KEY AUTOINCREMENT,
    wash_id INTEGER,
    phase VARCHAR(30),
    start_time DATETIME,
    end_time DATETIME,
    sensor_readings TEXT, -- JSON storage of sensor data
    FOREIGN KEY (wash_id) REFERENCES wash_history(id)
);
CREATE INDEX idx_cycle_log_wash ON cycle_log(wash_id);
\end{lstlisting}

\section{Module-wise Detailed Design}

\subsection{Vision Module}

The vision module is responsible for capturing and analyzing fabric images to determine fabric types.

\textbf{Components:}
\begin{itemize}
    \item Image Capture Controller: Manages camera initialization, capture timing, and image storage.
    \item Image Preprocessor: Applies filters, normalization, and augmentation to optimize images for inference.
    \item Fabric Classifier: Loads trained CNN model and performs inference on preprocessed images.
    \item Multi-Fabric Handler: Analyzes multiple items in a load and determines dominant or most delicate fabric.
\end{itemize}

\textbf{Algorithm:}
\begin{enumerate}
    \item Initialize camera with optimal settings (focus, exposure, white balance).
    \item Capture multiple images with different lighting angles.
    \item For each image:
        \begin{enumerate}
            \item Preprocess (resize to 224x224, normalize pixel values).
            \item Run inference through CNN model.
            \item Store confidence scores for each fabric class.
        \end{enumerate}
    \item Aggregate results across images (weighted average).
    \item If multiple items detected:
        \begin{enumerate}
            \item Segment image to identify individual clothing items.
            \item Classify each item separately.
            \item Determine wash strategy based on fabric composition.
        \end{enumerate}
    \item Return fabric type(s) and confidence scores.
\end{enumerate}

\subsection{Sensor Module}

The sensor module continuously monitors physical parameters throughout the wash cycle.

\textbf{Components:}
\begin{itemize}
    \item Weight Sensor Interface: Reads load cells to determine total laundry weight.
    \item Turbidity Sensor Interface: Monitors water clarity for soil level detection.
    \item Temperature Sensor Interface: Tracks water temperature for verification.
    \item Sensor Fusion Engine: Combines sensor data for comprehensive load characterization.
\end{itemize}

\textbf{Algorithm:}
\begin{enumerate}
    \item At cycle start:
        \begin{enumerate}
            \item Read load cells to determine dry weight.
            \item Classify load size (small/medium/large).
        \end{enumerate}
    \item During fill:
        \begin{enumerate}
            \item Monitor turbidity to assess soil level.
            \item Track temperature rise during heating.
        \end{enumerate}
    \item During wash:
        \begin{enumerate}
            \item Monitor turbidity changes to detect soil release.
            \item Monitor drum vibration for imbalance detection.
        \end{enumerate}
    \item During rinse:
        \begin{enumerate}
            \item Monitor turbidity to verify rinse effectiveness.
        \end{enumerate}
    \item Return all sensor data to control module.
\end{enumerate}

\subsection{Cycle Generation Module}

This module calculates optimal wash parameters based on fabric types, load weight, and soil level.

\textbf{Components:}
\begin{itemize}
    \item Parameter Calculator: Determines water volume, detergent quantity, and cycle duration.
    \item Resource Optimizer: Minimizes resource usage while maintaining cleaning effectiveness.
    \item Cycle Selector: Chooses appropriate wash cycle from knowledge base.
\end{itemize}

\textbf{Algorithm:}
\begin{enumerate}
    \item Input: fabric\_type(s), load\_weight, soil\_level.
    \item Query fabric\_knowledge database for base parameters.
    \item Calculate water volume: base\_water = fabric\_type.water\_factor × load\_weight; adjust for soil level (+10\% for heavy soil).
    \item Calculate detergent: base\_detergent = fabric\_type.detergent\_factor × load\_weight; adjust for soil level (+15\% for heavy soil); adjust for fabric delicacy (-20\% for delicates).
    \item Determine cycle parameters: temperature = fabric\_type.default\_temp; duration = base\_duration + soil\_adjustment; spin\_speed = fabric\_type.spin\_speed; rinse\_count = fabric\_type.delicate\_flag ? 3 : 2.
    \item Return complete cycle specification.
\end{enumerate}

\subsection{Control Module}

The control module executes the wash cycle by coordinating all actuators and monitoring progress.

\textbf{Components:}
\begin{itemize}
    \item State Machine: Manages cycle phases (fill, wash, rinse, spin, drain).
    \item Actuator Controller: Sends signals to valves, pump, dispenser, and motor.
    \item Safety Monitor: Detects errors and initiates safe shutdown if needed.
    \item Progress Tracker: Updates UI with current phase and time remaining.
\end{itemize}

\textbf{State Machine:}

\begin{figure}[H]
\centering
\fbox{\begin{minipage}{0.8\textwidth}
\centering
\textbf{Control State Machine}
\begin{itemize}
    \item IDLE $\rightarrow$ ANALYZE (on start)
    \item ANALYZE $\rightarrow$ FILL (cycle ready)
    \item FILL $\rightarrow$ HEAT (level OK)
    \item HEAT $\rightarrow$ WASH (temp OK)
    \item WASH $\rightarrow$ DRAIN (time up)
    \item DRAIN $\rightarrow$ SPIN (empty)
    \item SPIN $\rightarrow$ RINSE? if more rinses, go to FILL; else go to FINAL SPIN $\rightarrow$ COMPLETE
\end{itemize}
\end{minipage}}
\caption{Control State Machine}
\label{fig:state}
\end{figure}

\subsection{User Interface Module}

The UI module provides real-time information and control options to the user.

\textbf{Components:}
\begin{itemize}
    \item Display Controller: Manages touchscreen display content.
    \item Input Handler: Processes touch events and button presses.
    \item Status Renderer: Shows current cycle phase, time remaining, resource usage.
    \item Settings Manager: Stores and retrieves user preferences.
    \item Mobile API: Optional REST API for mobile app connectivity.
\end{itemize}

\textbf{Screen Layouts:}

\begin{figure}[H]
\centering
\begin{minipage}{0.6\textwidth}
\centering
\textbf{Main Screen Layout}
\begin{lstlisting}[basicstyle=\ttfamily\small, breaklines=true]
---------------------------------
AI Washing Machine         12:30
---------------------------------
Detected Fabrics:
  - Cotton 65%
  - Polyester 35%
Load Weight: 3.2 kg
Soil Level: Medium
Recommended Cycle:
Mixed Fabric
Water: 45L  Detergent: 35ml
Duration: 52 min
[  START  ]  [  CUSTOMIZE  ]
---------------------------------
\end{lstlisting}
\end{minipage}
\caption{Main Screen Layout}
\label{fig:ui}
\end{figure}

\section{Hardware-Software Interface}

\subsection{Communication Protocol}

The Raspberry Pi and Arduino communicate via serial USB with a simple command protocol:

\begin{table}[H]
\centering
\caption{Serial Communication Protocol}
\label{tab:protocol}
\begin{tabular}{|l|l|p{5.5cm}|}
\hline
\textbf{Command} & \textbf{Direction} & \textbf{Description} \\
\hline
GET\_WEIGHT & Pi $\rightarrow$ Arduino & Request load cell reading \\
WEIGHT:1234 & Arduino $\rightarrow$ Pi & Return weight in grams \\
GET\_TURB & Pi $\rightarrow$ Arduino & Request turbidity reading \\
TURB:345 & Arduino $\rightarrow$ Pi & Return turbidity value (0-1023) \\
GET\_TEMP & Pi $\rightarrow$ Arduino & Request temperature reading \\
TEMP:28.5 & Arduino $\rightarrow$ Pi & Return temperature in °C \\
VALVE\_HOT:ON & Pi $\rightarrow$ Arduino & Turn hot water valve on \\
VALVE\_HOT:OFF & Pi $\rightarrow$ Arduino & Turn hot water valve off \\
VALVE\_COLD:ON & Pi $\rightarrow$ Arduino & Turn cold water valve on \\
VALVE\_COLD:OFF & Pi $\rightarrow$ Arduino & Turn cold water valve off \\
DISPENSE:25 & Pi $\rightarrow$ Arduino & Dispense 25ml detergent \\
MOTOR:SPIN:800 & Pi $\rightarrow$ Arduino & Set motor to 800 RPM spin \\
MOTOR:AGITATE & Pi $\rightarrow$ Arduino & Start agitation cycle \\
MOTOR:STOP & Pi $\rightarrow$ Arduino & Stop motor \\
PUMP:ON & Pi $\rightarrow$ Arduino & Turn drain pump on \\
PUMP:OFF & Pi $\rightarrow$ Arduino & Turn drain pump off \\
STATUS & Pi $\rightarrow$ Arduino & Request system status \\
STATUS:OK & Arduino $\rightarrow$ Pi & System OK response \\
ERROR:code & Arduino $\rightarrow$ Pi & Error notification \\
\hline
\end{tabular}
\end{table}

\subsection{Arduino Control Code Structure}

\begin{lstlisting}[language=C++, caption={Arduino Main Loop Pseudocode}, label={lst:arduino}]
void loop() {
    if (Serial.available()) {
        String command = Serial.readStringUntil('\n');
        
        if (command.startsWith("GET_WEIGHT")) {
            float weight = readLoadCells();
            Serial.print("WEIGHT:");
            Serial.println(weight);
        }
        else if (command.startsWith("GET_TURB")) {
            int turb = readTurbidity();
            Serial.print("TURB:");
            Serial.println(turb);
        }
        else if (command.startsWith("VALVE_HOT:ON")) {
            digitalWrite(HOT_VALVE_PIN, HIGH);
            Serial.println("STATUS:OK");
        }
        else if (command.startsWith("VALVE_HOT:OFF")) {
            digitalWrite(HOT_VALVE_PIN, LOW);
            Serial.println("STATUS:OK");
        }
        // ... handle other commands
    }
    
    // Continuous monitoring
    checkSensors();
    delay(100);
}
\end{lstlisting}

\section{Safety and Error Handling}

\subsection{Safety Features}

\begin{table}[H]
\centering
\caption{Safety Features}
\label{tab:safety}
\small
\begin{tabular}{|l|p{3.2cm}|p{3.2cm}|p{3.8cm}|}
\hline
\textbf{Safety Feature} & \textbf{Description} & \textbf{Trigger} & \textbf{Action} \\
\hline
Water Leak Detection & Moisture sensors in base tray & Water detected & Close valves, activate alarm, drain pump \\
\hline
Overflow Prevention & Secondary water level sensor & Water above safe level & Close valves, activate drain pump \\
\hline
Temperature Limit & Redundant temperature monitoring & Water > 95°C & Stop heating, add cold water \\
\hline
Door Lock & Electromagnetic door lock & During operation & Prevent door opening \\
\hline
Imbalance Detection & Vibration sensor & Excessive vibration & Reduce spin speed, redistribute load \\
\hline
Motor Overcurrent & Current monitoring & Current > threshold & Stop motor, error display \\
\hline
Detergent Low & Level sensor in detergent tank & Detergent below minimum & User notification, prevent dispensing \\
\hline
\end{tabular}
\end{table}

\subsection{Error Codes and Handling}

\begin{table}[H]
\centering
\caption{Error Codes}
\label{tab:errors}
\small
\begin{tabular}{|l|l|p{4.5cm}|p{4cm}|}
\hline
\textbf{Error Code} & \textbf{Description} & \textbf{User Message} & \textbf{Recovery} \\
\hline
E001 & Camera communication failed & "Camera not detected. Check connection." & Retry, use default cycle \\
\hline
E002 & Fabric classification failed & "Unable to identify fabrics. Using standard cycle." & Default to cotton cycle \\
\hline
E003 & Water fill timeout & "Water fill taking too long. Check water supply." & Stop cycle, user intervention \\
\hline
E004 & Drain failure & "Water not draining. Check drain hose." & Stop cycle, user intervention \\
\hline
E005 & Heating element failure & "Heater not working. Cycle continuing with cold water." & Continue with cold wash \\
\hline
E006 & Motor driver error & "Motor error. Service required." & Safe shutdown \\
\hline
E007 & Load cell calibration error & "Weight sensor error. Using estimated load." & Estimate based on drum fill \\
\hline
E008 & Door not locked & "Door not securely closed." & Prevent start \\
\hline
\end{tabular}
\end{table}

\section{Performance Optimization}

\subsection{Image Processing Optimization}
\begin{itemize}
    \item \textbf{Resolution Scaling:} Capture at 1080p but downscale to 224x224 for inference.
    \item \textbf{Frame Skipping:} Capture every 2 seconds during analysis, not continuous video.
    \item \textbf{ROI Detection:} Identify clothing regions before full classification.
    \item \textbf{Model Quantization:} Convert model to TensorFlow Lite for faster inference.
    \item \textbf{Hardware Acceleration:} Use Raspberry Pi's GPU for neural network inference.
\end{itemize}

\subsection{Resource Calculation Optimization}
\begin{itemize}
    \item \textbf{Pre-computed Lookup Tables:} Store optimal parameters for common fabric combinations.
    \item \textbf{Caching:} Store recently used fabric parameters in memory.
    \item \textbf{Incremental Updates:} Adjust calculations based on previous cycle success.
    \item \textbf{Machine Learning Integration:} Use regression models for fine-tuned resource prediction.
\end{itemize}

\subsection{Real-time Performance Targets}

\begin{table}[H]
\centering
\caption{Performance Targets}
\label{tab:perf}
\begin{tabular}{|l|l|l|}
\hline
\textbf{Operation} & \textbf{Target Time} & \textbf{Measurement Method} \\
\hline
Fabric detection (single item) & $<$ 5 seconds & From image capture to classification \\
\hline
Full load analysis & $<$ 10 seconds & All items in drum \\
\hline
Cycle parameter calculation & $<$ 1 second & After fabric detection \\
\hline
Sensor read cycle & $<$ 100 ms & All sensors polled \\
\hline
Actuator response & $<$ 50 ms & Command to physical action \\
\hline
UI refresh rate & 30 fps & Screen updates \\
\hline
\end{tabular}
\end{table}

\section{Summary}

This chapter has presented the complete system design of the AI-Integrated Washing Machine, spanning the four-layer architecture, hardware design, software architecture, database schema, module specifications, safety features, and performance optimization strategies. The design decisions reflect consistent priorities: accurate fabric detection, precise resource optimization, reliable operation, and user-friendly interaction.

The hardware design integrates a camera system, multiple sensors, and precision actuators controlled by a Raspberry Pi and Arduino. The software architecture implements a modular approach with separate vision, sensor, cycle generation, control, and user interface modules. The database schema provides persistent storage for fabric knowledge, wash history, and user preferences.

Safety features ensure reliable operation under various conditions, while error handling provides clear user guidance when problems occur. Performance optimization strategies target real-time responsiveness suitable for household appliance operation.

The following chapter will detail the implementation of each module described here, translating these design specifications into working code and hardware assembly.